\documentclass[10pt,letterpaper]{article}
\usepackage[top=0.50in,bottom=0.52in,left=0.62in,right=0.62in,headheight=14pt,headsep=5pt,footskip=15pt]{geometry}
\usepackage{fontspec}
\setmainfont{Noto Serif}
\setsansfont{Noto Sans}
\usepackage[spanish,es-noshorthands]{babel}
\usepackage{microtype,ragged2e,setspace}
\usepackage{graphicx}
\usepackage{booktabs,tabularx,array,multirow,longtable}
\usepackage[table]{xcolor}
\usepackage{tikz}
\usetikzlibrary{arrows.meta,positioning,calc,shapes.geometric,babel}
\usepackage{amsmath,amssymb,rotating}
\usepackage{fancyhdr,hyperref,enumitem}
\usepackage{multicol}

% Industrial/VIU palette
\definecolor{autumn}{HTML}{E36A2E}
\definecolor{spicy}{HTML}{C8572A}
\definecolor{onyx}{HTML}{0C0C0C}
\definecolor{mustard}{HTML}{FCD757}
\definecolor{canary}{HTML}{8CB369}
\colorlet{orange}{autumn}
\colorlet{ink}{onyx}
\colorlet{muted}{onyx!58!white}
\colorlet{rule}{onyx!18!white}
\definecolor{soft}{HTML}{F5F7FA}
\definecolor{softblue}{HTML}{EAF2FD}
\definecolor{softgreen}{HTML}{EAF7EF}
\definecolor{softwarm}{HTML}{FDEFE8}
\definecolor{softplum}{HTML}{F2ECFB}
\definecolor{gamecol}{HTML}{2D6CC0}
\definecolor{controlcol}{HTML}{E36A2E}
\definecolor{searchcol}{HTML}{2E9B5F}
\definecolor{learncol}{HTML}{7B4DBA}
\definecolor{foundcol}{HTML}{6F7C8E}

\hypersetup{colorlinks=true,urlcolor=spicy,linkcolor=spicy,pdftitle={Marco teorico y estado del arte - multi-AMR},pdfauthor={Jorge Luis Mayorga Taborda}}
\pagestyle{fancy}\fancyhf{}
\fancyhead[L]{\sffamily\fontsize{7.3}{8}\selectfont Jorge Luis Mayorga Taborda}
\fancyhead[R]{\sffamily\fontsize{7.3}{8}\selectfont TFM MROB · Marco teórico y estado del arte}
\fancyfoot[C]{\sffamily\fontsize{8.2}{8.7}\selectfont\thepage}
\renewcommand{\headrulewidth}{0.25pt}
\setlength{\parindent}{0pt}\setlength{\parskip}{2.2pt}
\setstretch{1.05}\color{ink}\emergencystretch=1.4em
\setlist[itemize]{leftmargin=4.3mm,itemsep=.8pt,topsep=.7pt,parsep=0pt}
\setlist[enumerate]{leftmargin=4.7mm,itemsep=.6pt,topsep=.7pt,parsep=0pt}

\newcommand{\sect}[1]{\par\vspace{0pt}{\sffamily\fontsize{14.5}{15.5}\selectfont\bfseries\color{orange}#1\par}\vspace{1.2pt}\textcolor{rule}{\rule{\linewidth}{0.35pt}}\vspace{3pt}}
\newcommand{\subh}[1]{\par\vspace{2.2pt}{\sffamily\fontsize{9.3}{10.2}\selectfont\bfseries\color{orange}#1\par}\vspace{.5pt}}
\newcommand{\figcap}[2]{\par\vspace{1.0pt}{\setstretch{1}\fontsize{8.1}{8.9}\selectfont\textbf{#1.}~\textit{#2}\par}\vspace{1.0pt}}
\newcommand{\sourcecap}[1]{\par{\setstretch{1}\fontsize{7.0}{7.7}\selectfont\itshape\centering Fuente: #1\par}}
\newcommand{\note}[1]{\par{\setstretch{1}\fontsize{7.0}{7.8}\selectfont\color{muted}#1\par}}
\newcolumntype{L}[1]{>{\RaggedRight\arraybackslash}p{#1}}
\newcolumntype{C}[1]{>{\centering\arraybackslash}p{#1}}
\newcolumntype{Y}{>{\RaggedRight\arraybackslash}X}
\newcommand{\yes}{\tikz[baseline=-0.55ex]\fill[canary!85!onyx](0,0)circle(1.18mm);}
\newcommand{\parti}{\tikz[baseline=-0.55ex]{\begin{scope}\clip(0,0)circle(1.18mm);\fill[mustard](-1.3mm,-1.3mm) rectangle (0,1.3mm);\end{scope}\draw[onyx!55,line width=.28pt](0,0)circle(1.18mm);}}
\newcommand{\no}{\tikz[baseline=-0.55ex]\draw[onyx!48,line width=.32pt](0,0)circle(1.18mm);}
\newcommand{\mappointL}[7]{%
  \filldraw[fill=#7,draw=white,line width=.64pt] (#1,#2) circle (2.38mm);
  \node[text=white,font=\sffamily\bfseries\fontsize{5.55}{5.75}\selectfont] at (#1,#2){#5};
  \node[anchor=west,align=left,text=#7!70!black,font=\sffamily\bfseries\fontsize{5.35}{5.60}\selectfont] at ($(#1,#2)+(0.036,0)$) {#3\,\textsuperscript{[#6]}\\[-.45mm]\fontsize{4.68}{4.95}\selectfont #4};}
\newcommand{\mappointR}[7]{%
  \filldraw[fill=#7,draw=white,line width=.64pt] (#1,#2) circle (2.38mm);
  \node[text=white,font=\sffamily\bfseries\fontsize{5.55}{5.75}\selectfont] at (#1,#2){#5};
  \node[anchor=east,align=right,text=#7!70!black,font=\sffamily\bfseries\fontsize{5.35}{5.60}\selectfont] at ($(#1,#2)+(-0.036,0)$) {#3\,\textsuperscript{[#6]}\\[-.45mm]\fontsize{4.68}{4.95}\selectfont #4};}
% Contrato tipográfico: todo texto interno de figuras usa Noto Sans.
\newcommand{\chip}[2]{\tikz[baseline=-0.6ex]{\filldraw[fill=#1,draw=white,line width=.35pt] (0,0) circle (1.75mm);}\,\textcolor{ink}{\sffamily #2}}

\begin{document}

% ================= PAGE 1 =================
\sect{Marco teórico y estado del arte}
\fontsize{9.35}{10.15}\selectfont
La revisión examina cómo la literatura conecta cuatro decisiones que una misión de transporte ejecuta de manera continua: formar la coalición, comprobar su factibilidad física, mover la carga y recuperarse de un fallo. El análisis responde tres preguntas: qué mecanismos cubren cada decisión, qué evidencia respalda esa cobertura y dónde permanecen interfaces sin resolver.

Se trabajó con dos niveles de evidencia. El núcleo canónico reúne 59 trabajos (38 CORE, 15 ENABLING y 6 fundacionales o de revisión) para la comparación detallada; el corpus analítico de 244 documentos sustenta las tendencias posteriores. Los preprints y las versiones publicadas se enlazaron para no duplicar líneas de trabajo. La búsqueda y el \textit{snowballing} cubrieron MRTA, formación de coaliciones, transporte de carga compartida, contacto y wrench, redes imperfectas y recuperación. ``Distribuido'' se codificó por autoridad de decisión, información disponible, dependencia de líder o servidor y mecanismo de cierre. No se infirió solo a partir de la presencia de control local.
\normalsize

El primer resultado sitúa los trabajos del núcleo según autoridad y explicitud del mecanismo. La posición es descriptiva; no representa desempeño ni calidad.

\figcap{Figura 1}{Mapa metodológico agrupado del corpus académico. Horizontal: autoridad de decisión (descentralizado $\rightarrow$ centralizado). Vertical: explicitud del mecanismo (white-box $\rightarrow$ data-driven). El color identifica la familia primaria. Las posiciones son ordinales y se desplazan solo para evitar solapes.}
\begin{center}
\begin{tikzpicture}[x=80mm,y=54mm,font=\sffamily]
  \fill[softblue] (-1,-1) rectangle (0,0);
  \fill[softgreen] (0,-1) rectangle (1,0);
  \fill[softplum!55!white] (-1,0) rectangle (0,1);
  \fill[softwarm] (0,0) rectangle (1,1);
  \foreach \x in {-1,-.75,-.5,-.25,0,.25,.5,.75,1} \draw[rule,line width=.28pt] (\x,-1)--(\x,1);
  \foreach \y in {-1,-.75,-.5,-.25,0,.25,.5,.75,1} \draw[rule,line width=.28pt] (-1,\y)--(1,\y);
  \draw[-{Stealth[length=2.3mm]},ink,line width=.68pt] (-1.04,0)--(1.075,0);
  \draw[-{Stealth[length=2.3mm]},ink,line width=.68pt] (0,-1.04)--(0,1.09);
  \node[anchor=south,font=\sffamily\bfseries\fontsize{7.2}{7.5}\selectfont] at (0,1.10) {data-driven};
  \node[anchor=north,font=\sffamily\bfseries\fontsize{7.2}{7.5}\selectfont] at (0,-1.09) {white-box};
  \node[rotate=90,anchor=south,font=\sffamily\bfseries\fontsize{7}{7.3}\selectfont] at (-1.12,0) {descentralizado};
  \node[rotate=90,anchor=north,font=\sffamily\bfseries\fontsize{7}{7.3}\selectfont] at (1.12,0) {centralizado};
  \node[align=center,font=\sffamily\bfseries\fontsize{6.35}{6.65}\selectfont] at (-.60,.90) {distribuido\\[-.35mm]data-driven};
  \node[align=center,font=\sffamily\bfseries\fontsize{6.35}{6.65}\selectfont] at (.60,.90) {centralizado\\[-.35mm]data-driven};
  \node[align=center,font=\sffamily\bfseries\fontsize{6.35}{6.65}\selectfont] at (-.60,-.90) {distribuido\\[-.35mm]white-box};
  \node[align=center,font=\sffamily\bfseries\fontsize{6.35}{6.65}\selectfont] at (.60,-.90) {centralizado\\[-.35mm]white-box};

  % data-driven, distributed / mixed
  \mappointL{-0.82}{0.72}{Bezerra}{2025}{B}{20}{learncol}
  \mappointL{-0.64}{0.55}{Huo}{2025}{H}{19}{learncol}
  \mappointL{-0.46}{0.73}{Dai}{2024}{D}{17}{learncol}
  \mappointL{-0.78}{0.35}{Shibata-RL}{2023}{S}{16}{learncol}
  \mappointL{-0.43}{0.32}{Pandit}{2026}{P}{28}{learncol}
  \mappointR{-0.13}{0.50}{Eoh/Park}{2021}{E}{29}{learncol}

  % data-driven, centralized/mixed
  \mappointL{0.34}{0.64}{SADCHER}{2026}{S}{25}{learncol}
  \mappointL{0.62}{0.42}{CF-HMRTA}{2025}{C}{21}{searchcol}

  % white-box distributed - staggered
  \mappointL{-0.90}{-0.18}{Pereira}{2004}{P}{1}{controlcol}
  \mappointL{-0.78}{-0.34}{Parker/Tang}{2006}{PT}{2}{foundcol}
  \mappointL{-0.88}{-0.54}{Vig/Adams}{2006}{V}{3}{gamecol}
  \mappointL{-0.69}{-0.17}{CBBA}{2009}{C}{4}{gamecol}
  \mappointL{-0.64}{-0.39}{Service}{2011}{S}{5}{foundcol}
  \mappointL{-0.57}{-0.60}{Chen/Sun}{2012}{CS}{6}{gamecol}
  \mappointL{-0.45}{-0.22}{Jang}{2018}{J}{8}{gamecol}
  \mappointL{-0.39}{-0.47}{Dutta}{2019}{D}{10}{gamecol}
  \mappointL{-0.36}{-0.67}{Dutta}{2021}{D}{12}{gamecol}
  \mappointL{-0.14}{-0.28}{Ebel}{2024}{E}{18}{controlcol}
  \mappointR{-0.03}{-0.50}{Shibata}{2023}{S}{14}{controlcol}
  \mappointR{-0.04}{-0.71}{Fukao}{2025}{F}{22}{controlcol}
  \mappointL{-0.48}{-0.90}{GRAPE-S}{2026}{G}{24}{gamecol}

  % white-box central/mixed
  \mappointL{0.08}{-0.18}{Alonso-Mora}{2017}{AM}{7}{controlcol}
  \mappointL{0.20}{-0.40}{Machado}{2019}{M}{9}{controlcol}
  \mappointL{0.31}{-0.65}{Verma repl.}{2019}{V}{11}{controlcol}
  \mappointL{0.47}{-0.22}{Elaamery}{2021}{E}{13}{controlcol}
  \mappointL{0.59}{-0.44}{Gong}{2023}{G}{15}{controlcol}
  \mappointL{0.76}{-0.22}{Rizzo}{2020}{R}{30}{controlcol}
  \mappointR{0.89}{-0.40}{Loh}{2012}{L}{31}{controlcol}
  \mappointR{0.83}{-0.64}{Yufka}{2015}{Y}{32}{controlcol}
  \mappointR{0.61}{-0.73}{Muhammed}{2026}{M}{26}{controlcol}
  \mappointR{0.34}{-0.88}{Sahu}{2026}{S}{27}{controlcol}
  \mappointL{0.02}{-0.90}{Tian}{2025}{T}{23}{controlcol}

  \node[star,star points=5,star point ratio=2.1,fill=orange,draw=white,line width=.8pt,minimum size=8.5mm] at (-0.18,-0.78) {};
  \node[anchor=west,font=\sffamily\bfseries\fontsize{7.0}{7.4}\selectfont,align=left] at (-0.09,-0.75) {Este TFM\\integración};
\end{tikzpicture}
\end{center}
\vspace{-0.8mm}
\begin{center}\fontsize{8.35}{8.9}\selectfont
\chip{gamecol}{juegos/mercados}\hspace{5.5mm}
\chip{controlcol}{control/consenso}\hspace{5.5mm}
\chip{searchcol}{búsqueda/heurística}\hspace{5.5mm}
\chip{learncol}{política aprendida}\hspace{5.5mm}
\chip{foundcol}{fundacional}
\end{center}
\sourcecap{elaboración propia a partir de los trabajos [1]-[32] y la codificación del corpus académico.}
\note{El mapa separa con claridad la pertenencia lógica a una coalición de su ejecución física. Esa distinción ordena la comparación de capacidades de la página siguiente y evita atribuir a un algoritmo de asignación garantías de contacto, movimiento o recuperación que no evalúa.}

\vspace{0.4mm}
\begin{center}
\begin{tikzpicture}[x=1mm,y=1mm,font=\sffamily]
\foreach \x/\w/\bg/\tit/\n/\txt in {
0/45/softblue/Coalition formation/18/{matching · auctions · games},
47/45/softwarm/Shared-load transport/23/{control · wrench · caging},
94/45/softgreen/Network and safety/15/{consensus · CBF · recovery},
141/45/softplum/Foundational/6/{population games · GNE}}{
  \fill[\bg,rounded corners=.8pt](\x,0)rectangle +(\w,8.6);
  \draw[rule,rounded corners=.8pt,line width=.28pt](\x,0)rectangle +(\w,8.6);
  \node[anchor=west,font=\sffamily\bfseries\fontsize{5.5}{5.8}\selectfont] at (\x+2,5.9){\tit};
  \node[anchor=east,font=\sffamily\bfseries\fontsize{8.1}{8.1}\selectfont,text=orange] at (\x+\w-2,5.9){\n};
  \node[anchor=west,font=\sffamily\fontsize{4.45}{4.75}\selectfont,text=muted] at (\x+2,2.0){\txt};}
\end{tikzpicture}
\end{center}
\vspace{-0.5mm}
\figcap{Figura 1b}{Resumen funcional del corpus por familia temática.}
\sourcecap{elaboración propia a partir del dataset académico.}

\newpage
% ================= PAGE 2 =================
La primera comparación funcional aplica la misma regla a los trabajos más próximos al problema del TFM. Matching, subastas y juegos hedónicos cubren atomicidad, heterogeneidad y decisión local en combinaciones distintas [4]-[6],[8],[10],[12],[24]. Alonso-Mora, Ebel, Shibata, Fukao, Tian y Muhammed, por su parte, formulan y validan aspectos del transporte cooperativo distribuido o descentralizado [7],[14],[18],[22],[23],[26]. La separación entre ambos frentes impide formular la brecha como una simple ausencia de control distribuido de carga.

\figcap{Figura 2}{Matriz de capacidades de los trabajos más discriminantes. Verde = evidencia explícita; amarillo = parcial/condicionada; círculo vacío = no.}
\begin{center}
\sffamily
\renewcommand{\arraystretch}{1.45}\setlength{\tabcolsep}{2.25pt}\fontsize{7.55}{8.0}\selectfont
\begin{tabularx}{\linewidth}{L{29mm}*{8}{C{11.6mm}}C{10.2mm}C{9.8mm}}
\toprule
\textbf{Trabajo} & \rotatebox{90}{\textbf{heterog.}} & \rotatebox{90}{\textbf{coal. var.}} & \rotatebox{90}{\textbf{reclut. local}} & \rotatebox{90}{\textbf{shared-load}} & \rotatebox{90}{\textbf{cert. física}} & \rotatebox{90}{\textbf{exec. dist.}} & \rotatebox{90}{\textbf{recovery}} & \rotatebox{90}{\textbf{replacement}} & \textbf{Ev.} & \textbf{C}\\
\midrule
Jang 2018 [8] & \no & \yes & \yes & \no & \no & \yes & \no & \no & T+S & C4\\
Dutta 2021 [12] & \yes & \yes & \yes & \no & \no & \yes & \no & \no & T+S & C4\\
GRAPE-S 2026 [24] & \yes & \yes & \parti & \no & \no & \yes & \no & \no & T+S & C3--4\\
CF-HMRTA 2025 [21] & \yes & \yes & \no & \no & \no & \no & \no & \no & T+S & C0--1\\
Huo 2025 [19] & \yes & \yes & \yes & \no & \no & \yes & \parti & \no & T+S & C3--4\\
Bezerra 2025 [20] & \yes & \yes & \yes & \no & \no & \yes & \parti & \no & S & C4\\
Alonso-Mora 2017 [7] & \no & \parti & \no & \yes & \yes & \parti & \parti & \no & T+S+P & C0--1\\
Ebel 2024 [18] & \no & \no & \no & \yes & \yes & \yes & \no & \no & T+S+P & C3\\
Shibata 2023 [14] & \no & \parti & \no & \yes & \yes & \yes & \parti & \parti & S+P & C4\\
Muhammed 2026 [26] & \no & \no & \no & \yes & \yes & \yes & \no & \no & T+S+P & C3\\
Verma repl. 2019 [11] & \no & \parti & \no & \yes & \parti & \no & \yes & \yes & S+P & C0--1\\
Sahu 2026 [27] & \no & \parti & \no & \yes & \parti & \parti & \yes & \yes & S+P & C1--2\\
\bottomrule
\end{tabularx}
\end{center}
\sourcecap{elaboración propia a partir de [7],[8],[11],[12],[14],[18]-[21],[24],[26],[27].}

La lectura conjunta conserva tres distinciones. Los juegos poblacionales y los GNE pueden producir una masa continua, de modo que el cierre entero debe formar parte del método. La desigualdad $\sum_i c_i\ge d_k$ tampoco certifica el transporte: se requiere $w_k^{req}\in\mathcal W(C_k)$ bajo geometría, fricción y límites de actuador. Finalmente, la escala debe separar $N_{sim}$, $N_{hw}$, robots por coalición y cargas simultáneas. Un experimento con mil robots simulados no equivale a una validación física de la misma escala.

\begin{center}
\begin{tikzpicture}[x=1mm,y=1mm,font=\sffamily\fontsize{6.25}{6.65}\selectfont]
\tikzset{bx/.style={rounded corners=1pt,draw=rule,line width=.35pt,minimum width=41.5mm,minimum height=17mm,align=center}}
\node[bx,fill=softblue](s)at(21,10){\textbf{Supported cargo}\\carga soportada\\wrench y límites};
\node[bx,fill=softwarm](r)at(65,10){\textbf{Rigid attachment}\\formación rígida\\distribución de fuerza};
\node[bx,fill=softgreen](p)at(109,10){\textbf{Pushing / caging}\\contacto unilateral\\geometría no-prehensil};
\node[bx,fill=softplum](g)at(153,10){\textbf{Grasping / arms}\\grasp matrix\\force closure};
\foreach \a/\b in {s/r,r/p,p/g}{\draw[-{Stealth[length=1.55mm]},muted,line width=.42pt](\a)--(\b);} 
\end{tikzpicture}
\end{center}
\figcap{Figura 2b}{Modalidades físicas que no deben mezclarse bajo la etiqueta genérica \textit{cooperative transport}.}
\sourcecap{síntesis conceptual basada en [1],[7],[9],[13]-[15],[18],[22],[23],[26],[28],[30]-[32].}

\newpage
% ================= PAGE 3 =================
Al añadir la continuidad de la misión, la cobertura de componentes aislados deja de ser suficiente. La cuestión decisiva es si esos componentes coinciden en una misma cadena operativa y si la autoridad permanece local cuando la misión cambia de etapa. La Figura 3 contrasta ambas dimensiones sobre el núcleo curado.

\figcap{Figura 3}{Síntesis estructural en dos vistas complementarias: \textbf{(a)} cobertura de intersecciones del corpus CORE y \textbf{(b)} presupuesto de centralización por etapa.}
\begin{center}
\begin{minipage}[t]{0.47\linewidth}
\centering
{\sffamily\bfseries\fontsize{7.4}{7.9}\selectfont (a) Cobertura del corpus CORE\par}
\vspace{1.5mm}
\begin{tikzpicture}[x=1.72mm,y=6.45mm,font=\sffamily\fontsize{5.95}{6.35}\selectfont]
\foreach \lab/\v/\y in {Shared-load transport/23/7,Recruit/18/6,Certificación física/13/5,Shared-load + distributed/10/4,Heterog. + coalición variable/9/3,Shared-load + replacement/2/2,Heterog. + shared-load/0/1,Physical + distributed + replacement/0/0}{
  \draw[rule,line width=.25pt] (0,\y) -- (24.8,\y);
  \fill[orange] (0,\y-.23) rectangle (\v,\y+.23);
  \node[anchor=east,align=right,text width=31mm] at (-1.0,\y) {\lab};
  \node[anchor=west,font=\sffamily\bfseries] at (\v+.45,\y) {\v};}
\draw[-{Stealth[length=1.45mm]},ink,line width=.4pt](0,-.65)--(25.3,-.65);
\node[anchor=north]at(12.2,-.82){trabajos CORE};
\end{tikzpicture}
\vspace{0.3mm}
\sourcecap{conteos del dataset académico curado; no constituyen prevalencias universales.}
\end{minipage}\hfill
\begin{minipage}[t]{0.49\linewidth}
\centering
{\sffamily\bfseries\fontsize{7.4}{7.9}\selectfont (b) Presupuesto de centralización por etapa\par}
\vspace{1.5mm}
\begin{tikzpicture}[x=0.90mm,y=0.90mm,font=\sffamily\fontsize{5.9}{6.25}\selectfont]
\draw[ink,line width=.45pt] (28,8)--(148,8);
\foreach \x/\lab in {28/local,68/mixto,108/global,148/central} {\draw[ink,line width=.3pt](\x,6.7)--(\x,9.3);\node[anchor=north]at(\x,5.7){\lab};}
\foreach \y/\lab/\x in {52/Recruit/45,41/Certify/76,30/Dock/48,19/Transport/44,8/Recover/79}{
  \node[anchor=east,font=\sffamily\bfseries]at(23,\y){\lab};
  \draw[rule,line width=.35pt](28,\y)--(148,\y);
  \fill[soft](28,\y-3)rectangle(148,\y+3);
  \fill[orange,rounded corners=.7pt](28,\y-3)rectangle(\x,\y+3);
  \draw[orange,line width=.45pt](\x,\y-3.8)--(\x,\y+3.8);} 
\end{tikzpicture}
\vspace{0.3mm}
\sourcecap{codificación de autoridad e información de los trabajos más cercanos [7],[11],[14],[18]-[21],[24],[26],[27].}
\end{minipage}
\end{center}
\fontsize{8.25}{8.95}\selectfont
Las dos vistas convergen. El corpus cubre transporte de carga compartida, reclutamiento y certificación física por separado, pero las intersecciones se reducen al añadir heterogeneidad, ejecución distribuida y sustitución. Al mismo tiempo, la autoridad residual cambia entre reclutamiento, acoplamiento, transporte y recuperación. La discontinuidad se encuentra en el paso entre etapas, no en la ausencia absoluta de cada ingrediente.
\normalsize

\rowcolors{2}{white}{soft}\fontsize{7.3}{7.8}\selectfont
\begin{tabularx}{\linewidth}{L{25mm} C{10mm} C{10mm} C{12mm} Y}
\toprule
\textbf{Trabajo} & \textbf{Share} & \textbf{Repl.} & \textbf{Local rec.} & \textbf{Diferencia decisiva}\\\midrule
Verma 2019 [11] & \yes & \yes & \no & scheduling global; hubs/ruta conocidos\\
Sahu 2026 [27] & \yes & \yes & \no & arquitectura master/slave\\
Shibata 2023 [14] & \yes & \parti & \no & equipo variable, no selección desde una flota mayor\\
Ebel 2024 [18] & \yes & \no & \no & equipo fijo\\
Bezerra 2025 [20] & \no & \no & \yes & no shared-load físico\\
GRAPE-S 2026 [24] & \no & \no & \parti & servicio lógico, no carga compartida\\
\bottomrule
\end{tabularx}

\vspace{1.2mm}
\fontsize{8.25}{8.95}\selectfont
De esta comparación se obtiene una delimitación precisa. En el corpus revisado no se identificó una arquitectura que reúna, dentro de una misma cadena operativa, reclutamiento local de una coalición heterogénea y variable desde una flota mayor, certificación física de la carga compartida y sustitución de un miembro durante el transporte, sin mantener una autoridad global permanente. Esta formulación restringe la contribución a la interfaz que permanece abierta y evita reclamar como nuevos los componentes ya resueltos.
\normalsize
\sourcecap{síntesis comparativa de [7],[11],[14],[18]-[21],[24],[26],[27].}

\vspace{1.0mm}
\fontsize{8.55}{9.25}\selectfont
Esta delimitación fija qué debe buscarse en el corpus ampliado. El análisis temporal comprueba si los métodos convergen hacia esa integración o si el crecimiento reciente continúa distribuyéndose entre capas separadas.
\normalsize

\fontsize{8.55}{9.25}\selectfont
La misma delimitación cambia la unidad de comparación. El análisis deja de clasificar cada trabajo mediante una única etiqueta global. Un sistema puede decidir localmente durante el transporte y depender de una asignación central para formar el equipo; también puede detectar un fallo de manera distribuida y delegar la selección del reemplazo en un servidor. La autoridad, la información disponible y la condición de cierre deben codificarse por etapa.

Esta regla permite comparar resultados que, de otro modo, parecerían equivalentes. SP1 debe entregar una coalición entera, roles y capacidad verificable; SP2 debe devolver margen físico y causas de fallo; SP3 debe incorporar la huella de la carga y el coste de congestión. La evolución del corpus puede examinarse entonces bajo una misma lectura y esos intercambios pueden convertirse en hipótesis y métricas del TFM.
\normalsize

\end{document}
